% ============================================================
% KDMILE 2026 — STAB Framework Paper — Versão PT-BR
% Deadline: 08/jun/2026 | Limite: 8 páginas
% Submissão: https://jems3.sbc.org.br/kdmile2026
% ============================================================

\documentclass{kdmile}

\usepackage{graphicx,url}
\usepackage[T1]{fontenc}
\usepackage[utf8]{inputenc}
\usepackage{booktabs}
\usepackage{multirow}
\usepackage{listings}
\usepackage{float}

% kdmile.cls é baseado em ACM Trans — define resumo e abstract manualmente
\newenvironment{abstracten}{%
  \vspace{4pt}\begin{center}\normalsize\bfseries Abstract\end{center}%
  \noindent\small\ignorespaces
}{\normalsize\vspace{4pt}}

\newenvironment{resumo}{%
  \vspace{4pt}\begin{center}\normalsize\bfseries Resumo\end{center}%
  \noindent\small\ignorespaces
}{\normalsize\vspace{4pt}}

\providecommand{\address}[1]{}
\renewcommand{\institutename}{}

% ── Cabeçalho ───────────────────────────────────────────────
\markboth{Nogueira et al.}{STAB: A Complexity-Ladder Framework for Evaluating Local LLMs}

% ── Título (obrigatório em inglês pelo call for papers) ──────
\title{STAB: A Complexity-Ladder Framework for Evaluating Local LLMs as Autonomous Agents}

% ── Autores ─────────────────────────────────────────────────
\author{
  João Conrado V. Nogueira\\
  {\small Pesquisador Independente, Macapá, AP, Brasil}
}

% ============================================================
\begin{document}
\maketitle
% Corrige rodapé: cls hardcoda KDMiLe 2025; sobrescreve após \maketitle
\gdef\@runningfoot{\@journalName, KDMiLe 2026.}
\gdef\@firstfoot{\@journalName, KDMiLe 2026.}
% ============================================================

% ── Abstract (EN — obrigatório pelo call for papers) ────────
\begin{abstracten}
Running LLMs as autonomous agents on local hardware is harder than it looks.
%
This paper reports STAB: a four-level complexity ladder (S1$\rightarrow$S4)
combined with a 12-hour temporal stability battery, built to expose
failure modes in on-premise tool-calling agents.
%
We tested eight models from four families (Qwen, Gemma, Llama, Mistral)
on a single NVIDIA RTX~3060 (12~GB VRAM).
%
The central finding is that most failures trace back to the instruction and
contextual decision layer—\allowbreak not to raw model capacity.
%
Two phenomena stood out: \emph{tool-skipping} triggered by descriptive
filenames (H12), and the orthogonality between format compliance and semantic
accuracy (H10).
%
\texttt{gemma4:e4b} was the only model to pass all four levels without a
single failure.

\textbf{Keywords:} Local LLMs, Autonomous Agents, Tool Calling, Instruction
Following, Small Language Models, On-Premise Benchmark.
\end{abstracten}

% ── Resumo (PT) ─────────────────────────────────────────────
\begin{resumo}
O uso de Large Language Models (LLMs) em sistemas agênticos locais enfrenta
desafios críticos de consistência e confiabilidade, especialmente em hardware
de consumo.
%
Este trabalho apresenta o protocolo STAB, uma metodologia de avaliação em escada
de complexidade (S1$\rightarrow$S4) combinada a um protocolo de estabilidade
temporal para diagnosticar falhas em agentes on-premise.
%
Avaliamos oito modelos de quatro famílias (Qwen, Gemma, Llama e Mistral) em
servidor com GPU NVIDIA RTX~3060 (12~GB VRAM).
%
Os resultados revelam que as falhas agênticas originam-se predominantemente na
camada de instrução e decisão contextual, e não na capacidade bruta dos modelos.
%
Identificamos o fenômeno de \emph{tool-skipping} por nomeação semântica (H12) e
demonstramos que conformidade de formato e acurácia semântica são dimensões
ortogonais (H10).
%
O modelo \texttt{gemma4:e4b} destacou-se como referência de maior
confiabilidade, mantendo desempenho perfeito em todos os níveis de complexidade.

\textbf{Palavras-chave:} LLMs Locais, Agentes Autônomos, Tool Calling,
Instruction Following, Small Language Models, Benchmark On-Premise.
\end{resumo}

% ============================================================
\section{Introdução}
\label{sec:intro}

Fazer com que um LLM local invoque ferramentas de forma confiável—\allowbreak e não apenas
responda perguntas—\allowbreak é mais difícil do que parece.
Modelos que se saem bem em benchmarks padrão frequentemente falham em contextos
agênticos: pulam invocações de ferramentas, alucinam resultados ou retornam JSON
correto encapsulado em markdown que o script chamador não consegue parsear.
Avaliações como BFCL e AgentBench não expõem esses problemas porque operam com
\textit{compute} praticamente ilimitado e permitem re-tentativas—\allowbreak condições que
não existem em hardware restrito com $T=0$ e execução one-shot.

Este artigo apresenta o STAB para preencher essa lacuna. As principais contribuições são:

\begin{itemize}
  \item \textbf{Framework STAB:} Metodologia de quatro níveis (S1$\rightarrow$S4)
        que avalia desde o uso básico de ferramentas até o planejamento autônomo
        de múltiplas etapas, validada por um protocolo de estabilidade temporal.

  \item \textbf{Taxonomia de Falhas:} Classificação de erros agênticos em cinco
        camadas: protocolo, decisão, formato, grounding e instrução.

  \item \textbf{H12 — Tool-Skipping Semântico:} Identificação empírica da
        omissão de invocação de ferramenta induzida por nomes de arquivo
        semanticamente descritivos.

  \item \textbf{H10 — Ortogonalidade:} Evidência de que conformidade de formato
        e acurácia semântica são dimensões de avaliação independentes.
\end{itemize}

O artigo está organizado da seguinte forma: a Seção~\ref{sec:related} revisa os
trabalhos relacionados; a Seção~\ref{sec:method} descreve a metodologia STAB;
a Seção~\ref{sec:results} apresenta resultados e análises; a
Seção~\ref{sec:discussion} discute os achados; e a Seção~\ref{sec:conclusion}
conclui o trabalho.

% ============================================================
\section{Trabalhos Relacionados}
\label{sec:related}

\textbf{Benchmarks de agentes e uso de ferramentas.}
O Berkeley Function Calling Leaderboard (BFCL)~\cite{patil2025} é o benchmark
de referência para avaliação de ferramentas, com foco predominante em modelos
em nuvem e LLMs de grande porte.
O AgentBench~\cite{liu2024} propõe avaliação em múltiplos ambientes em oito
domínios, identificando o seguimento de instruções como principal obstáculo à
autonomia.
O SLM-Bench~\cite{pham2025} introduz métricas de impacto ambiental e
configurações de hardware controladas, aproximando-se de nossa metodologia de
avaliação local.

\textbf{Modos de falha e seguimento de instruções.}
Em~\cite{cemri2025} é proposta uma taxonomia de 14 modos de falha em
sistemas multiagente, destacando a desobediência a especificações.
O desempenho de SLMs para function calling em dispositivos de borda é mapeado
em~\cite{kavathekar2025}, enquanto~\cite{jhandi2026} demonstra que SLMs com até
350M parâmetros podem superar modelos maiores no uso de ferramentas com ajuste
fino direcionado. Uma queda na conformidade com instruções em modelos menores é
documentada por meio do dataset KCIF em~\cite{murthy2024}, e pipelines de
treinamento para cenários corporativos que exigem uso determinístico de
ferramentas são explorados em~\cite{zeng2024}.

\textbf{Padrões fundacionais de agentes com ferramentas.}
O paradigma ReAct~\cite{yao2023} intercala raciocínio e ação em um mesmo loop de
geração, base conceitual da maioria dos agentes de tool-calling atuais. O
Toolformer~\cite{schick2023} demonstrou que LLMs podem aprender a decidir
\emph{quando} invocar uma ferramenta de forma auto-supervisionada; nenhum dos
dois avalia essa robustez sob restrição de hardware ou execuções repetidas. Uma
linha ortogonal à engenharia de prompt ataca a conformidade de formato na
camada de decodificação: \cite{geng2023} restringe o espaço de tokens gerados a
saídas estruturalmente válidas via gramática formal—\allowbreak técnica discutida como
mitigação ao anti-padrão AP2 na Seção~\ref{sec:discussion}.

\textbf{Diferenciação.}
O STAB tem um posicionamento diferente: foca em hardware on-premise, introduz
uma bateria de estabilidade temporal e sonda explicitamente vieses contextuais
como nomeação semântica—\allowbreak nenhum desses aspectos aparece nos benchmarks acima.

% ============================================================
\section{Metodologia}
\label{sec:method}

\subsection{Configuração Experimental}

\begin{table}[ht]
\centering
\caption{Ambiente experimental.}
\label{tab:setup}
\begin{tabular}{ll}
\toprule
\textbf{Componente} & \textbf{Especificação} \\
\midrule
CPU     & Intel Core i5-10400F (6c/12t) \\
RAM     & 32\,GB DDR4 \\
GPU     & NVIDIA RTX 3060 (12\,GB VRAM) \\
SO      & Ubuntu 24.04 \\
Runtime & Ollama (REST \texttt{/api/chat}), $T=0$, \texttt{stream=false} \\
\bottomrule
\end{tabular}
\end{table}

\subsection{Framework STAB — Escada de Complexidade}

O framework STAB organiza a avaliação em quatro níveis de complexidade, cada
um exigindo maior capacidade cognitiva e operacional do agente:

\begin{table}[ht]
\centering
\caption{Escada de complexidade STAB.}
\label{tab:stab}
\begin{tabular}{clllcc}
\toprule
\textbf{Nível} & \textbf{Nome} & \textbf{Ferramenta} & \textbf{Objetivo} & \textbf{K} & \textbf{Métrica} \\
\midrule
S1 & Protocolo     & \texttt{run\_bash}  & Invocação + retorno de ferramenta & 5 & TOOL-ACC \\
S2 & Extração      & \texttt{run\_bash}  & Saída JSON estruturada            & 5 & JSON-ACC \\
S3 & Análise de Log& \texttt{read\_file} & Classificação semântica           & 4 & CLASS-ACC \\
S4 & Multi-log     & \texttt{read\_file} & Triagem e ranking multi-etapa     & 3 & TRIAGE-ACC \\
\bottomrule
\end{tabular}
\end{table}

O framework utiliza fixtures sintéticas em vez de logs de produção para
garantir reprodutibilidade e controle de variáveis. Cada fixture consiste em
um arquivo de entrada (.log) e um arquivo JSON de referência (.expected.json).
A referência emprega vocabulário controlado para identificação de problemas
(e.g., \texttt{high\_cpu\_usage}, \texttt{database\_unavailable}),
transformando a avaliação aberta em classificação e extração rigorosas.

Cada teste segue um protocolo de execução em dois turnos via API REST do Ollama
(\texttt{/api/chat}). No Turno~1, o modelo recebe o prompt e o JSON Schema das
ferramentas, produzindo uma \texttt{tool\_call} estruturada. O script executa a
operação no host e captura o \texttt{stdout}. No Turno~2, esse resultado é
retornado ao modelo, que gera a resposta JSON final. O determinismo é garantido
configurando a temperatura como zero e \texttt{stream=false}.

\subsection{Modelos Avaliados}

\begin{table}[ht]
\centering
\caption{Modelos avaliados.}
\label{tab:models}
\begin{tabular}{llrll}
\toprule
\textbf{Modelo} & \textbf{Família} & \textbf{VRAM (MiB)} & \textbf{Tipo} & \textbf{Status Final} \\
\midrule
\texttt{gemma4:e4b}             & Gemma4    & 9{,}821 & Zero-shot  & \textbf{Referência} \\
\texttt{qwen3.5:9b}             & Qwen3.5   & 7{,}953 & Zero-shot  & Aprovado (S2B) \\
\texttt{qwen3:14b}              & Qwen3     & 9{,}437 & Zero-shot  & Aprovado (S1--S3) \\
\texttt{llama3.1:8b}            & Llama3.1  & 4{,}868 & Zero-shot  & Aprovado (Parcial) \\
\texttt{llama3-groq-tool-use:8b}& Llama3    & 4{,}754 & Fine-tuned & Aprovado (Parcial) \\
\texttt{mistral:7b}             & Mistral   & 4{,}458 & Zero-shot  & Descartado (H14) \\
\texttt{phi4-mini}              & Phi4      & 2{,}500 & Zero-shot  & Descartado (H13) \\
\texttt{qwen3.5:2b}             & Qwen3.5   & 3{,}829 & Zero-shot  & Descartado (H16) \\
\bottomrule
\end{tabular}
\end{table}

\subsection{Métricas}

A avaliação emprega uma escala ordinal de 0 a 4, onde 4 representa uma execução
perfeita (invocação correta da ferramenta, saída JSON válida e conteúdo semântico
preciso) e 0 representa falha completa.

\textbf{Definição de falha.} \emph{Falha}: execução que não atinge escore
máximo (4) em alguma métrica aplicável ao nível. \emph{Falha crítica}—\allowbreak descarte
do modelo, interrompe a progressão—\allowbreak é um escore 0 em métrica bloqueante do nível
(TOOL-ACC em S1; JSON-ACC em S2), já que o nível seguinte pressupõe o anterior
como pré-condição. Falhas não críticas (escore 1--3) mantêm o modelo na
progressão e são registradas na taxonomia (Seção~\ref{sec:method}) e na
Tabela~\ref{tab:results}.

As métricas estão agrupadas em três
categorias: métricas de consistência (CON, TOOL-CON) avaliam a estabilidade
entre execuções; métricas de instrução e acurácia (INST, TOOL-ACC, JSON-ACC,
CLASS-ACC, PARSE-ACC, CHAIN-ACC, TRIAGE-ACC) medem a correção em cada nível
STAB; e o escore composto STAB reflete a estabilidade temporal ao longo do
protocolo completo de bateria de 12 horas.

\begin{table}[ht]
\centering
\caption{Métricas de avaliação (escala ordinal 0--4).}
\label{tab:metrics}
\begin{tabular}{ll}
\toprule
\textbf{Métrica} & \textbf{Descrição} \\
\midrule
CON       & Consistência: estabilidade estrutural e decisória entre execuções \\
TOOL-CON  & Consistência de ferramenta: invocações válidas repetidas \\
INST      & Aderência a instrução: conformidade de formato e restrições \\
TOOL-ACC  & Correção funcional: ferramenta correta, execução bem-sucedida \\
JSON-ACC  & Acurácia JSON: validade do schema e precisão de valores \\
CLASS-ACC & Acurácia de classificação: severidade global correta \\
PARSE-ACC & Acurácia de extração: host, problema e timestamp precisos \\
CHAIN-ACC & Acurácia de encadeamento: número correto de chamadas (e.g. 3/3) \\
TRIAGE-ACC& Acurácia de triagem: ranking de prioridade e regra de escalada \\
STAB      & Estabilidade temporal: manutenção de escores em baterias de 12\,h \\
\bottomrule
\end{tabular}
\end{table}

\subsection{Protocolo de Estabilidade Temporal}
\label{sec:stab-protocol}

A métrica STAB (Tabela~\ref{tab:metrics}) exige um protocolo operacional
específico. Cada modelo aprovado executa o conjunto completo de cenários STAB
duas vezes—\allowbreak Bateria~1 e Bateria~2—\allowbreak com intervalo mínimo de 12h, hardware,
prompts, temperatura e ordem de cenários idênticos entre elas. O intervalo
isola variância técnica temporal (aquecimento da GPU, cache do SO, drift de
estado do Ollama entre reinicializações) da variância intrínseca ao modelo. O
escore STAB de um cenário é a diferença absoluta entre Bateria~1 e Bateria~2
nessa métrica; STAB=4 indica escore idêntico nas duas. Divergências são
investigadas inspecionando os \emph{traces} de raciocínio interno para
distinguir instabilidade real de variação estrutural (ver \texttt{qwen3:14b}
na Seção~\ref{sec:results}).

\subsection{Taxonomia de Falhas}

As falhas foram classificadas em cinco camadas com base em sua causa raiz.
Falhas na camada de Protocolo envolvem violações estruturais do formato de
chamada de ferramenta. Falhas na camada de Decisão indicam erros de raciocínio
contextual, como inferir resultados sem executar a ferramenta necessária. Falhas
na camada de Formato produzem saídas sintaticamente corretas, mas não conformes.
Falhas de Grounding geram respostas estruturalmente válidas, porém factualmente
fabricadas. Falhas na camada de Instrução decorrem de interpretação ambígua de
ações. Cada falha recebe um código único: códigos com prefixo H indicam achados
orientados por hipóteses confirmados por experimentos controlados, enquanto
códigos AP marcam padrões anômalos identificados durante execuções padrão.

\begin{table}[ht]
\centering
\caption{Taxonomia de falhas identificadas.}
\label{tab:failures}
\begin{tabular}{lllll}
\toprule
\textbf{Camada} & \textbf{Tipo} & \textbf{Cód.} & \textbf{Exemplo Concreto} & \textbf{Modelo} \\
\midrule
Protocolo   & Alucinação de Protocolo & H13 & JSON em \texttt{content} em vez de \texttt{tool\_calls} & phi4-mini \\
Protocolo   & Echo-tool               & H16 & Entrega JSON via \texttt{run\_bash echo}                & qwen3.5:2b \\
Decisão     & Tool-Skipping Semântico & H12 & Leitura omitida; severidade inferida do nome do arquivo & qwen3:14b \\
Formato     & Markdown Wrapper        & AP2 & JSON correto encapsulado em \texttt{```json}            & qwen3.5:9b \\
Grounding   & Grounding Silencioso    & H14 & JSON perfeito com dados de hardware fabricados          & mistral:7b \\
Instrução   & Ambiguidade de Ação     & AP1 & Usou \texttt{cat} em vez de executar o script           & qwen3.5:9b \\
\bottomrule
\end{tabular}
\end{table}

% ============================================================
\section{Resultados}
\label{sec:results}

\subsection{Comparação Geral}

Três modelos não passaram pelos níveis iniciais: \texttt{phi4-mini} (H13, S1),
\texttt{qwen3.5:2b} (H16, S2) e \texttt{mistral:7b} (H14, S2)—\allowbreak cada um por um
motivo estrutural diferente, detalhados a seguir.

O \texttt{phi4-mini} falha na camada de \emph{protocolo}: alternou entre
serializar o JSON da ferramenta no campo \texttt{content} (em vez de
\texttt{tool\_calls}) e ignorar a ferramenta, alucinando um resultado plausível
no texto. Por não reconhecer a estrutura esperada da API independentemente da
instrução, nenhuma correção via prompt foi tentada—\allowbreak essa classe de falha não é
candidata natural à engenharia de prompt.

Para \texttt{mistral:7b} e \texttt{qwen3.5:2b}, correções \emph{foram}
tentadas, com resultados instrutivos. O \texttt{mistral:7b} recebeu instrução
reforçada em S2B ("use \texttt{run\_bash} com o comando exato")—\allowbreak continuou sem
invocar a ferramenta, mas o prompt refinado eliminou o wrapper de markdown que
antes acionava o validador, produzindo um \textbf{falso positivo}: JSON
perfeito (13/13 campos) com dados fabricados (timestamp de três anos no
passado, GPU e RAM inexistentes no servidor). Isso motiva a distinção entre
H14 e as demais falhas de protocolo (Tabela~\ref{tab:failures}): o validador de
schema sozinho não capturaria a falha. O \texttt{qwen3.5:2b} recebeu instrução
igualmente direta, testada via chamada direta à API—\allowbreak o wrapper de markdown
persistiu, sugerindo causa no template de chat ou treinamento, não ambiguidade
corrigível em inferência.

Dos cinco que prosseguiram, o \texttt{gemma4:e4b} foi o único a passar em todos
os níveis sem falha (S1$\rightarrow$S4; ver Tabela~\ref{tab:results}).
A queda mais acentuada ocorreu no S4: modelos que haviam tratado S1--S3 sem
problemas colapsaram nas tarefas de triagem multi-etapa, apontando para um teto
na camada de decisão que tarefas de etapa única simplesmente não revelam.

\begin{table}[ht]
\centering
\caption{Escores STAB por modelo e nível (0--4; \textbf{negrito} = execução perfeita).}
\label{tab:results}
\begin{tabular}{lccccccl}
\toprule
\textbf{Modelo} & \textbf{S1} & \textbf{S2} & \textbf{S3} & \textbf{S4-A} & \textbf{S4-B} & \textbf{S4-C} & \textbf{S1→S4} \\
\midrule
\texttt{gemma4:e4b}              & \textbf{4} & \textbf{4} & \textbf{4} & \textbf{4} & \textbf{4} & \textbf{4} & \textbf{Sim} \\
\texttt{qwen3.5:9b}              & 4 & 4 (S2B) & 2$^a$ & 4 & 2 & 2 & Não \\
\texttt{qwen3:14b}               & 4 & 4 & 4 & 0 & 0 & 0 & Não \\
\texttt{llama3.1:8b}             & 4 & 4 & 2 & 4 & 2 & 3 & Não \\
\texttt{llama3-groq-tool-use:8b} & 4 & 4 & 4 & 4 & 2 & 2 & Não \\
\texttt{qwen3.5:2b}              & 4 & desc. H16 & --- & --- & --- & --- & Não \\
\texttt{mistral:7b}              & 4 & desc. H14 & --- & --- & --- & --- & Não \\
\texttt{phi4-mini}               & desc. H13 & --- & --- & --- & --- & --- & Não \\
\bottomrule
\end{tabular}
\begin{minipage}{\linewidth}
\footnotesize
$^a$\,\texttt{qwen3.5:9b} atingiu CLASS-ACC=4 no S3-mixed, mas CON=INST=2 devido ao
markdown wrapper persistente, ilustrando a ortogonalidade das dimensões (H10).
\end{minipage}
\end{table}

\subsection{Principais Achados}

\textbf{H12 — Tool-Skipping Semântico.}
O modelo \texttt{qwen3:14b} omitiu a leitura de arquivos em cenários S4 quando
os nomes eram semanticamente descritivos (e.g., \texttt{critical-001.log}),
inferindo a severidade sem executar a ferramenta. No experimento controlado H12
com nomes de arquivo opacos, o modelo invocou corretamente a ferramenta,
confirmando a hipótese de viés contextual. \textbf{Mitigação:} utilizar nomes
de arquivo opacos ou instruções explícitas de leitura obrigatória.

\textbf{H10 — Ortogonalidade de Formato e Acurácia Semântica.}
Conformidade de formato e acurácia semântica mostraram-se dimensões
independentes. Um modelo pode produzir saída semanticamente perfeita encapsulada
em formato inválido (markdown), ou um JSON estruturalmente válido contendo dados
alucinados. Isso foi exemplificado pelo \texttt{qwen3.5:9b}, que manteve alta
acurácia semântica enquanto falhava nas restrições de formatação.

\textbf{gemma4:e4b como Modelo de Referência.}
O \texttt{gemma4:e4b} foi o único modelo a atingir escores perfeitos em todos os
níveis (S1$\rightarrow$S4) com STAB=4 e throughput de 73{,}0\,tok/s. Seu
raciocínio interno (\emph{thinking}) foi ativado seletivamente para cenários de
desambiguação de alta complexidade (S4-C, como desempate baseado em timestamp),
indicando escalonamento cognitivo adaptativo.

\subsection{Throughput e Latência}

\begin{figure}[ht]
  \centering
  \includegraphics[width=\linewidth]{plot-throughput-kdmile.png}
  \caption{Throughput médio por modelo e cenário (tok/s), Bateria~1.
           Execução~1 excluída (cold start). Variação $<0{,}2\%$ entre baterias
           confirma inferência limitada pelo hardware.}
  \label{fig:throughput}
\end{figure}

A variação entre baterias ficou abaixo de 0,2\% para todos os modelos,
confirmando que o throughput é limitado pelo hardware, não pelo prompt.
A Execução~1 de cada bateria foi excluída das médias (penalidade de cold start).
Três clusters de desempenho emergiram: \texttt{gemma4:e4b} e a família Llama no
topo, \texttt{qwen3.5:9b} no meio e \texttt{qwen3:14b} na base—\allowbreak sua latência no
S2 inflada por um extenso trace de raciocínio interno.

\subsection{Estabilidade Temporal}

Todos os cinco modelos aprovados completaram a Bateria~2 sem degradação:
decisões agênticas e acurácia semântica mantiveram STAB=4 ao longo do protocolo.
A única anomalia foi o \texttt{qwen3:14b}, cujo volume de raciocínio interno
variou 7\% entre baterias—\allowbreak mas a saída final foi idêntica em ambos os casos,
de modo que isso representa uma peculiaridade estrutural do chain-of-thought,
não instabilidade real.

% ============================================================
\section{Discussão}
\label{sec:discussion}

\textbf{Falhas são instrucionais, não baseadas em capacidade.}
Dois experimentos corretivos neste estudo contam a mesma história.
Uma instrução de formato explícita recuperou o escore S2B do \texttt{qwen3.5:9b}
de 2 para 4. Uma diretiva de leitura obrigatória eliminou completamente o
tool-skipping do \texttt{qwen3:14b}.
Nenhuma das correções exigiu trocar o modelo ou fazer retreinamento—\allowbreak apenas o
prompt.
O tamanho do modelo é portanto um proxy fraco para confiabilidade agêntica:
modelos acima de 8B claramente possuem a lógica subjacente; o que os quebra é
o desalinhamento de contexto, e isso é corrigível na camada de instrução.

\textbf{Fine-tuning vs.\ zero-shot em ambientes on-premise.}
O confronto entre \texttt{llama3-groq-\allowbreak tool-\allowbreak use:8b} (fine-tuned) e
\texttt{llama3.1:8b} (zero-shot) produziu um resultado dividido.
No S3-mixed, que exige aplicar regras de precedência binária, o modelo
fine-tuned venceu com folga (Groq=4 vs.\ Llama=2)—\allowbreak especialização funciona para
conformidade estrutural em etapa única.
Mas no S4-B e S4-C, ambos falharam.
Isso coloca uma pergunta concreta: o fine-tuning para uso de ferramentas melhora
a aderência estrutural sem tocar no raciocínio multi-etapa necessário para
decisões encadeadas? Os dados aqui sugerem que sim—\allowbreak e isso é uma limitação real
dos pipelines de ajuste fino atuais.

\textbf{Recomendações práticas.}
Em hardware com 12\,GB de VRAM, \texttt{gemma4:e4b} é a escolha óbvia:
consistência perfeita, 73\,tok/s, sem necessidade de engenharia de prompt.
Abaixo de 6\,GB, Llama3.1 funciona—\allowbreak mas precisa de uma camada fina de
pós-processamento para tratar falhas de formato.
O resultado H12 tem também um ângulo de segurança: nomes de arquivo descritivos
permitem que o modelo infira respostas e pule a chamada de ferramenta, o que
significa que a ferramenta nunca é de fato executada.
Usar identificadores opacos, ou acrescentar \texttt{[LEITURA OBRIGATÓRIA]} à
instrução, fecha essa brecha.

\textbf{Decodificação com restrição gramatical como mitigação agnóstica ao
modelo.}
O AP2 (wrapper de markdown) e a resistência à instrução do
\texttt{qwen3.5:2b} (Seção~\ref{sec:results}) são falhas de \emph{formato}, não
de conteúdo—\allowbreak o modelo "sabe" a resposta, mas a emite de um jeito que o parser
não aceita. Engenharia de prompt resolve esse problema só parcialmente
(recuperou o \texttt{qwen3.5:9b}, não o \texttt{qwen3.5:2b}). Geração restrita
por gramática formal~\cite{geng2023}—\allowbreak via GBNF (\texttt{llama.cpp}), Outlines ou
\texttt{format=json} do Ollama—\allowbreak ataca a camada de \emph{decodificação} em vez da
de entrada, tornando a violação de formato impossível independentemente da
cooperação do modelo. É candidata natural exatamente para os casos resistentes
a instrução direta, como o \texttt{qwen3.5:2b}. O STAB não testou essa técnica
(Seção~\ref{sec:conclusion}), mas os dados aqui apontam o caso de uso onde ela
seria mais valiosa.

\textbf{Limitações.}
O estudo rodou em uma única configuração de hardware e utilizou apenas fixtures
sintéticas. São restrições reais, mas que também tornaram possível a
reprodutibilidade total e o isolamento de variáveis—\allowbreak um trade-off deliberado.
O número de execuções por nível também é reduzido (K=5/5/4/3; Tabela~\ref{tab:stab}),
limitando o poder estatístico e favorecendo a leitura dos resultados como
diagnóstico qualitativo, não como taxa de erro com significância estatística.
O resultado mais relevante operacionalmente a levar adiante é o H14 (grounding
silencioso no Mistral): saída estruturalmente perfeita que é factualmente errada
não será capturada por um parser JSON.
Validadores semânticos externos não são opcionais em produção—\allowbreak são a única
proteção contra esse modo de falha específico.

% ============================================================
\section{Conclusão}
\label{sec:conclusion}

A confiabilidade de agentes locais é um problema de três variáveis: robustez do
modelo, clareza de instrução e neutralidade de contexto.
As três precisam funcionar simultaneamente—\allowbreak a falha em qualquer uma delas quebra
a cadeia.
O \texttt{gemma4:e4b} foi o único modelo neste estudo a manter as três em
equilíbrio em todos os níveis, mostrando que arquiteturas MoE quantizadas
cruzaram um limiar prático para trabalho agêntico estruturado em hardware de
consumidor.

As falhas nos modelos de 8B--14B foram, quase sem exceção, corrigíveis apenas
com engenharia de prompt—\allowbreak sem retreinamento.
Esse é um resultado encorajador para quem constrói agentes sob restrições de
hardware.

O que este estudo não cobriu: logs de produção reais, perfis de hardware
diversificados e decodificação com restrição gramatical como alternativa ao
enforcement de formato baseado em prompt.
Fixtures sintéticas (Seção~\ref{sec:method}) garantiram reprodutibilidade, mas
não capturam o ruído de logs reais—\allowbreak timestamps malformados, encoding
inconsistente, entradas truncadas, múltiplos problemas concorrentes. O nível S4
(triagem multi-log), já o de queda mais acentuada mesmo com fixtures limpas
(Seção~\ref{sec:results}), é o candidato mais provável a se degradar ainda mais
sob esse ruído—\allowbreak validação direta para trabalho futuro.
A bateria temporal de 12 horas também deixa em aberto se deployments mais longos
expõem padrões de degradação que não foram visíveis aqui.
Essas são as extensões naturais—\allowbreak não por convenção, mas porque os dados deste
estudo levantam essas questões específicas.

% ============================================================
\bibliographystyle{kdmile}
\bibliography{references}

\end{document}
